\chapter{双原子分子的振动与转动}\label{ch:diatomic-vib-rot}

前面的 Born--Oppenheimer 思想表明：在选定一条电子势能面之后，双原子分子的低能内部运动首先是两个原子核在该势能曲线上的量子运动。本章从两核 Hamiltonian 开始，不预先假设“振动”和“转动”彼此独立，而是先精确分离质心，再得到只含核间距与取向的相对运动方程。谐振子、刚性转子以及离心畸变都将作为这一母方程的受控近似依次出现。

\section{一般多原子分子的质量加权正常模}

Born--Oppenheimer 势能面 $E_n(\mathbf R)$ 建立以后，振动理论首先是一个一般的多自由度二次型问题，而双原子分子只是只有一个内部振动坐标的特例。设平衡构型为 $\mathbf R^{(0)}$，笛卡尔位移记为 $\delta R_{A\alpha}$。在平衡点附近
\[
E_n(\mathbf R)=E_n(\mathbf R^{(0)})
+\frac12\sum_{A\alpha,B\beta}
K_{A\alpha,B\beta}\delta R_{A\alpha}\delta R_{B\beta}+O(\delta R^3),
\]
其中 Hessian $K$ 的元素量纲为力常数。定义质量加权坐标 $q_{A\alpha}=\sqrt{M_A}\,\delta R_{A\alpha}$ 与质量加权 Hessian
\begin{equation}
F_{A\alpha,B\beta}
=\frac{K_{A\alpha,B\beta}}{\sqrt{M_AM_B}},
\qquad
\sum_{B\beta}F_{A\alpha,B\beta}L_{B\beta,s}=\omega_s^2L_{A\alpha,s}.
\label{eq:mass-weighted-hessian}
\end{equation}
去除整体平移与转动零模以后，剩余正本征值给出分子正常模频率 $\omega_s$。非线性 $N$ 原子分子有 $3N-6$ 个振动模，线性分子有 $3N-5$ 个。

\begin{derivation}[为什么 Hessian 必然包含平移与转动零模]
\emph{第一步：平移不变性。}孤立分子势能 $V(\{\mathbf R_A\})$ 只依赖粒子间相对几何，因此对整体平移 $\mathbf R_A\mapsto\mathbf R_A+\mathbf a$ 不变：
\begin{equation*}
V(\{\mathbf R_A^{(0)}+\mathbf a+\delta\mathbf R_A\})
=V(\{\mathbf R_A^{(0)}+\delta\mathbf R_A\}).
\end{equation*}
把两侧都在平衡点按 $\delta\mathbf R$ 展开。一阶项给出 $\sum_{A\alpha}\partial_{A\alpha}V\,a_\alpha=0$，它在平衡点自动成立（合力为零）；比较二阶项得，对任意平移生成向量 $t_{A\alpha}=a_\alpha$，
\begin{equation*}
\sum_{A\alpha,B\beta}K_{A\alpha,B\beta}\,t_{A\alpha}\,t_{B\beta}=0.
\end{equation*}
$K$ 对称且上式对任意 $\mathbf a$ 成立。取 $\mathbf a$ 沿单一坐标方向并利用对称双线性型恒为零推得
\begin{equation}
\sum_{B\beta}K_{A\alpha,B\beta}\,t_{B\beta}=0,
\qquad\text{即}\qquad K\,t=0.
\label{eq:hessian-translation-zero}
\end{equation}
三个独立平移方向 $\mathbf a=\mathbf e_x,\mathbf e_y,\mathbf e_z$ 给出 $K$ 的三个线性无关零模。物理上，零本征值意味着沿该方向移动不改变势能到二阶——整体平移不拉伸任何键，自然无恢复力。

\emph{第二步：旋转不变性。}无外场分子势能对整体旋转 $R$ 不变，$V(\{R\mathbf R_A\})=V(\{\mathbf R_A\})$。取无穷小旋转 $R=I+\boldsymbol\theta\times$，位移 $\delta\mathbf R_A=\boldsymbol\theta\times\mathbf R_A^{(0)}$，即 $t_{A\alpha}=(\boldsymbol\theta\times\mathbf R_A^{(0)})_\alpha$。与平移完全相同的二阶比较给出
\begin{equation*}
\sum_{B\beta}K_{A\alpha,B\beta}(\boldsymbol\theta\times\mathbf R_B^{(0)})_\beta=0,
\end{equation*}
对任意 $\boldsymbol\theta$ 成立，故三个无穷小转动方向也是 $K$ 的零模。若全部原子共线，绕分子轴的转动 $\boldsymbol\theta_\parallel\times\mathbf R_A^{(0)}=0$ 不产生任何位移，该方向退化，只剩两条独立转动零模。转动零模的显式构造：取 $\boldsymbol\theta=\mathbf e_x$，则 $t_{A}^{(x)}=(0,-R_{Az}^{(0)},R_{Ay}^{(0)})$；类似给出 $t_A^{(y)}$、$t_A^{(z)}$。这三个向量在质量加权后 $\tilde t_A^{(k)}=t_A^{(k)}/\sqrt{M_A}$ 一般与平移零模不正交，需要 Gram--Schmidt 处理。

\emph{零模计数。}非线性 $N$ 原子分子的 $3N$ 个笛卡尔自由度分解为：3 个整体平移（零模）+ 3 个整体转动（零模）+ $(3N-6)$ 个振动。线性分子绕轴的转动不产生位移，故转动零模只有 2 个，振动自由度增加为 $3N-5$。这一计数与分子点群的结构一致：$3N$ 维表示在平移表示 $T$、转动表示 $R$ 和振动表示 $\Gamma_{\rm vib}$ 上的分解 $3N=T\oplus R\oplus\Gamma_{\rm vib}$，$\Gamma_{\rm vib}$ 的维数正是 $3N-6$ 或 $3N-5$。

\emph{第三步：质量加权 Hessian 的零空间。}质量加权 Hessian $F_{A\alpha,B\beta}=K_{A\alpha,B\beta}/\sqrt{M_AM_B}$ 即 $F=M^{-1/2}KM^{-1/2}$。若 $Kt=0$，则 $F(M^{1/2}t)=M^{-1/2}Kt=0$；$M^{1/2}$ 可逆，所以 $F$ 与 $K$ 零空间维数相同。去掉零模后剩余正本征值个数即振动自由度：
\begin{equation*}
3N-3-3=3N-6\;(\text{非线性}),\qquad
3N-3-2=3N-5\;(\text{线性}).
\end{equation*}
质量加权坐标 $q_{A\alpha}=\sqrt{M_A}\delta R_{A\alpha}$ 下的动能是 $\frac12\sum_{A\alpha}\dot q_{A\alpha}^2$，不再含质量；本征向量 $L_{A\alpha,s}$ 直接给出正常模在质量加权坐标中的位移方向，转换回笛卡尔位移需除以 $\sqrt{M_A}$。

\emph{第四步：零模的物理意义与投影。}Hessian 的零模不是数值病态，而是分子对称性的精确代数后果。在实际计算中，必须把 $F$ 限制到零模补空间才能求振动频率。标准做法是构造 $3N\times(3N-6)$（或 $3N-5$）的投影矩阵 $P$，其列向量为质量加权的平移和转动零模的正交补基，然后对角化 $P^T F P$。具体步骤：(i) 构造 6 个（或 5 个）零模向量 $\{t_i\}$，平移零模 $t_{A\alpha}^{(k)}=\delta_{\alpha k}/\sqrt{M_A}$（质量加权归一化），转动零模 $t_{A\alpha}^{(k)}=(\mathbf e_k\times\mathbf R_A^{(0)})_\alpha/\sqrt{M_A}$；(ii) Gram--Schmidt 正交化 $\{t_i\}$；(iii) 构造投影 $P=I-\sum_i|t_i\rangle\langle t_i|$；(iv) 对角化 $PFP$ 的非零本征值。若误把零模当作振动模，会出现虚假的零频率"振动"，这是初学者常犯的错误。

\emph{零模与分子点群。}零模的对称性由分子点群 $G$ 的平移表示 $T$ 和转动表示 $R$ 给出。振动表示 $\Gamma_{\rm vib}=3N-T-R$ 可约化为 $G$ 的不可约表示直和，每个不可约表示出现的次数由 $\chi_{\rm vib}(g)=\chi_{3N}(g)-\chi_T(g)-\chi_R(g)$ 的特征标公式决定。这一分解预测红外活性（$\Gamma_{\rm vib}$ 含 $T$ 的不可约分量）和拉曼活性（含二阶对称张量分量），无需显式计算 Hessian。

\emph{简并模的对称性破缺。}若 Hessian 有简并正本征值，对应正常模构成分子点群的不可约表示。Jahn--Teller 效应表明：非线性分子的轨道简并电子态会使核构型畸变以破缺简并，对应 Hessian 在高对称构型处的负本征值。零模分析因此不仅给出振动频率，还检测平衡构型的稳定性：负本征值意味着该构型是过渡态而非稳定极小。

\emph{过渡态的零模与反应路径。}在化学反应中，过渡态（saddle point）的 Hessian 有且仅有一个负本征值（沿反应坐标），其余正本征值对应垂直于反应路径的振动模。过渡态理论的反应速率 $k=k_BT/h\cdot\ee^{-E^\ddagger/(k_BT)}$ 中，$E^\ddagger$ 是过渡态与反应物的能量差，配分函数比正是来自去除反应坐标零模（此处是"虚频"模）后的振动配分函数。这一连接把分子振动理论与化学动力学统一起来。

\emph{举例：水分子的振动模。}水分子（$N=3$，非线性）有 $3N-6=3$ 个振动模。质量加权 Hessian 去除 6 个零模（3 平移 + 3 转动）后，剩余 3 个正本征值对应：对称伸缩 $\nu_1\approx3657\,$cm$^{-1}$（两 O--H 键同相伸缩）、弯曲 $\nu_2\approx1595\,$cm$^{-1}$（H--O--H 键角变化）、反对称伸缩 $\nu_3\approx3756\,$cm$^{-1}$（两 O--H 键反相伸缩）。三个模的对称性分别属于 $C_{2v}$ 点群的 $A_1$、$A_1$、$B_2$ 不可约表示，红外活性由 $\partial\mu/\partial Q_s\ne0$ 判定（$\mu$ 是分子偶极矩），三个模均红外活性。这一分析完全建立在 Hessian 本征值问题之上，零模的精确去除是关键前提。

\emph{举例：线性 CO$_2$ 分子的振动模。}CO$_2$（$N=3$，线性）有 $3N-5=4$ 个振动模：对称伸缩 $\nu_1\approx1388\,$cm$^{-1}$（O--C--O 同相，$D_{\infty h}$ 的 $\Sigma_g^+$，红外非活性因偶极矩恒为零）、二重简并弯曲 $\nu_2\approx667\,$cm$^{-1}$（$\Pi_u$，二重简并来自两个正交弯曲平面，红外活性）、反对称伸缩 $\nu_3\approx2349\,$cm$^{-1}$（$\Sigma_u^+$，红外活性）。线性分子比非线性少一个振动模但多一个转动模（绕分子轴的转动退化），总数仍为 $3N-5+2=3N-3$ 内部自由度。

\emph{更一般结论：内坐标与 Wilson 矩阵。}对一般分子，可直接用内坐标（键长、键角、二面角）$s_i$ 构造 Wilson $G$ 矩阵和 $F$ 矩阵（力常数矩阵），振动频率由 $|GF-\lambda I|=0$ 给出。$G$ 矩阵自动编码了质量加权内坐标的动能，避免了显式去除平移转动零模的步骤。这是 Wilson、Decius、Cross 的 GF 矩阵方法，与质量加权 Hessian 本征值问题等价，但在大分子中数值更稳定。零模的群论根源——分子构型空间在平移转动群作用下的商结构 $Q=(\mathbb R^{3N}-\Delta)/E(3)$——是分子振动几何学的核心：振动光谱研究的不是 $\mathbb R^{3N}$ 上的 Hessian，而是商空间 $Q$ 上的约化 Hessian，零模正是商投影的核。

\emph{更一般结论：反应路径上的 Hessian 演化。}沿内禀反应坐标（IRC）跟踪 Hessian 本征值，可检测从稳定极小到过渡态的拓扑变化。IRC 上的 Hessian 分解把 $3N-6$ 维振动空间分为沿反应路径的"软"模和垂直的"硬"模，后者在反应过程中近似绝热跟随，前者决定反应动力学。这一分析把振动零模从静态平衡点推广到动态反应路径，也是过渡态搜索算法（如 eigenvector-following）的数学基础。

\emph{数值实现注记。}实际量子化学程序中，质量加权 Hessian 的对角化需用对称正定求解器，零模通过本征值阈值（典型 $10^{-6}\,$a.u.）识别并投影去除。
\end{derivation}

\section{双原子分子的精确两体约化}
接下来考察质心与相对坐标的精确分离。

设两个原子核质量分别为 \(m_A,m_B\)，实验室坐标为 \(\mathbf R_A,\mathbf R_B\)。在一条给定的绝热势能曲线 \(U(R)\) 上，核 Hamiltonian 为

\begin{equation}
H_N
=-\frac{\hbar^2}{2m_A}\nabla_A^2
-\frac{\hbar^2}{2m_B}\nabla_B^2
+U(R),
\qquad
R=|\mathbf R_A-\mathbf R_B|.
\label{eq:diatomic-nuclear-hamiltonian}
\end{equation}

定义总质量、质心坐标和相对坐标

\[
M=m_A+m_B,
\qquad
\mathbf R_{\rm cm}
=\frac{m_A\mathbf R_A+m_B\mathbf R_B}{M},
\qquad
\mathbf R=\mathbf R_A-\mathbf R_B.
\]

由链式法则，梯度变换为

\[
\nabla_A
=\frac{m_A}{M}\nabla_{\rm cm}+\nabla_R,
\qquad
\nabla_B
=\frac{m_B}{M}\nabla_{\rm cm}-\nabla_R.
\]

代入动能项，交叉项严格抵消：

\begin{align*}
\frac{\nabla_A^2}{m_A}+\frac{\nabla_B^2}{m_B}
&=\left(\frac{m_A}{M^2}+\frac{m_B}{M^2}\right)\nabla_{\rm cm}^2
+\left(\frac1{m_A}+\frac1{m_B}\right)\nabla_R^2\\
&\quad
+\frac{2}{M}\nabla_{\rm cm}\!\cdot\!\nabla_R
-\frac{2}{M}\nabla_{\rm cm}\!\cdot\!\nabla_R\\
&=\frac1M\nabla_{\rm cm}^2+\frac1\mu\nabla_R^2,
\end{align*}

其中

\begin{equation}
\boxed{
\mu=\frac{m_Am_B}{m_A+m_B}
}
\label{eq:diatomic-reduced-mass}
\end{equation}

是约化质量。因此 \eqnrefs{eq:diatomic-nuclear-hamiltonian} 精确分裂成

\begin{equation}
H_N
=-\frac{\hbar^2}{2M}\nabla_{\rm cm}^2
+H_{\rm int},
\qquad
H_{\rm int}
=-\frac{\hbar^2}{2\mu}\nabla_R^2+U(R).
\label{eq:diatomic-internal-hamiltonian}
\end{equation}

质心自由平移与内部光谱完全解耦，所以后文只研究 \(H_{\rm int}\)。这一步没有使用振动小振幅、刚性键长或慢转动等近似。

接下来考察球坐标分离与精确径向方程。

相对坐标 \(\mathbf R\) 同时包含键长 \(R\) 和分子轴取向 \((\theta,\phi)\)。球坐标中的 Laplace 算符可以写成

\begin{equation}
\nabla_R^2
=\frac1{R^2}\frac{\dd}{\dd R}
\left(R^2\frac{\dd}{\dd R}\right)
-\frac{\hat{\mathbf J}^{\,2}}{\hbar^2R^2},
\label{eq:radial-laplacian-angular}
\end{equation}

这里 \(\hat{\mathbf J}\) 是核相对运动的轨道角动量算符。因为 \(U(R)\) 是中心势，\([H_{\rm int},\mathbf J^2]=[H_{\rm int},J_Z]=0\)，可以取

\[
\Psi_{JM}(R,\theta,\phi)
=F_J(R)Y_{JM}(\theta,\phi),
\qquad
\mathbf J^2Y_{JM}=\hbar^2J(J+1)Y_{JM}.
\]

把它代回 Schr\"odinger 方程得到

\[
-\frac{\hbar^2}{2\mu}
\left[
\frac1{R^2}\frac{\dd}{\dd R}
\left(R^2\frac{\dd F_J}{\dd R}\right)
-\frac{J(J+1)}{R^2}F_J
\right]
+U(R)F_J
=EF_J.
\]

为了去掉径向一阶导数，定义约化径向波函数

\[
u_J(R)=R F_J(R).
\]

直接计算

\[
\frac1{R^2}\frac{\dd}{\dd R}
\left(R^2\frac{\dd}{\dd R}\frac{u}{R}\right)
=\frac1R\frac{\dd^2u}{\dd R^2},
\]

于是得到一维形式的精确径向方程

\begin{equation}
\boxed{
-\frac{\hbar^2}{2\mu}\frac{\dd^2u_J}{\dd R^2}
+
\left[
U(R)+\frac{\hbar^2J(J+1)}{2\mu R^2}
\right]u_J
=E u_J.
}
\label{eq:diatomic-radial-equation}
\end{equation}

因此所谓“振动--转动耦合”并不是后加的一项：它从一开始就包含在离心势

\[
U_{\rm cent}(R)=\frac{\hbar^2J(J+1)}{2\mu R^2}
\]

中。把振动和转动完全分开，只是 \eqnrefs{eq:diatomic-radial-equation} 的某个尺度极限。

\section{从一般势能曲线到 Dunham 展开}
接下来考察平衡位置附近的系统展开。

设 \(R_e\) 是电子势能曲线的稳定极小值，满足

\[
U'(R_e)=0,
\qquad
k\equiv U''(R_e)>0.
\]

令 \(q=R-R_e\)，Taylor 展开为

\begin{equation}
U(R_e+q)
=U_e+\frac12kq^2
+\frac1{3!}k_3q^3
+\frac1{4!}k_4q^4+\cdots,
\label{eq:diatomic-potential-expansion}
\end{equation}

其中 \(k_n=U^{(n)}(R_e)\)。最低非零项给出角频率

\begin{equation}
\omega_e=\sqrt{\frac{k}{\mu}},
\qquad
\ell_v\equiv\sqrt{\frac{\hbar}{\mu\omega_e}}.
\label{eq:diatomic-vibrational-scale}
\end{equation}

\(\ell_v\) 是振动波函数的典型空间宽度。谐近似的控制条件不是抽象的“\(q\) 很小”，而是低振动态主要取样的区域 \(|q|\sim\ell_v\) 内高阶项相对于二次项足够小，例如

\[
\epsilon_3
\sim\frac{|k_3|\ell_v}{k}\ll1,
\qquad
\epsilon_4
\sim\frac{|k_4|\ell_v^2}{k}\ll1.
\]

忽略 \(k_3,k_4,\ldots\) 后得到标准谐振子

\[
H_{\rm vib}^{(0)}
=-\frac{\hbar^2}{2\mu}\frac{\dd^2}{\dd q^2}
+\frac12\mu\omega_e^2q^2.
\]

其本征能量与归一化本征函数为

\begin{equation}
\boxed{
E_v^{\rm vib}
=\hbar\omega_e\left(v+\frac12\right),
\qquad
v=0,1,2,\ldots
}
\label{eq:diatomic-harmonic-levels}
\end{equation}

和

\[
\psi_v(q)
=\left(\frac{\alpha_v}{\pi}\right)^{1/4}
\frac{H_v(\sqrt{\alpha_v}q)}{\sqrt{2^vv!}}
\exp\!\left(-\frac12\alpha_vq^2\right),
\qquad
\alpha_v=\frac{\mu\omega_e}{\hbar}.
\]

这里的 \(v\) 是振动量子数。随着 \(v\) 增大，波函数向更远离 \(R_e\) 的区域延伸，高阶势能项必然变得重要，因此谐近似的失效首先发生在高振动态，而不是在所有能级上同时发生。

接下来考察刚性转子是径向方程的冻结极限。

若振动宽度满足

\[
\epsilon_R\equiv\frac{\ell_v}{R_e}\ll1,
\]

并且研究的 \(J\) 不大到足以显著拉长键长，就可以在离心项中先取 \(R\simeq R_e\)。定义转动惯量

\[
I_e=\mu R_e^2
\]

以及能量量纲的转动常数

\begin{equation}
B_e\equiv\frac{\hbar^2}{2I_e}
=\frac{\hbar^2}{2\mu R_e^2},
\label{eq:diatomic-rotational-constant}
\end{equation}

则刚性转子能级为

\begin{equation}
\boxed{
E_J^{\rm rot}=B_eJ(J+1),
\qquad
J=0,1,2,\ldots
}
\label{eq:diatomic-rigid-rotor}
\end{equation}

每个 \(J\) 在无外场时具有 \(2J+1\) 个 \(M\) 简并态。若谱学采用波数 \(\tilde\nu\) 表示能量，常定义

\[
\widetilde B_e=\frac{B_e}{hc}
=\frac{\hbar}{4\pi c I_e}.
\]

这里要区分两个不同的近似：\eqnrefs{eq:diatomic-harmonic-levels} 来自势能面的局部二次展开，而 \eqnrefs{eq:diatomic-rigid-rotor} 来自把离心项中的径向自由度冻结在 \(R_e\)。二者都从同一个精确径向方程出发，但控制它们的尺度并不完全相同。

接下来考察转振耦合与离心畸变。

真正的分子转动越快，离心势就越倾向于增大平均键长。这个效应可以直接从 \eqnrefs{eq:diatomic-radial-equation} 的有效势中计算，而不必把离心畸变常数 \(D\) 当成经验参数。

令

\[
L_J^2\equiv\hbar^2J(J+1).
\]

在低振动量子数和中等 \(J\) 下，对 \(q/R_e\) 展开

\[
\frac1{(R_e+q)^2}
=\frac1{R_e^2}
\left(1-\frac{2q}{R_e}+\frac{3q^2}{R_e^2}+\cdots\right).
\]

只保留产生最低阶 \(J^4\) 能量修正所需的线性项，得到

\begin{align*}
U_{\rm eff}(q)
&\simeq
U_e+\frac12kq^2
+\frac{L_J^2}{2\mu R_e^2}
-\frac{L_J^2}{\mu R_e^3}q.
\end{align*}

把关于 \(q\) 的两项配方：

\begin{align*}
\frac12kq^2-\frac{L_J^2}{\mu R_e^3}q
&=\frac12k(q-q_J)^2
-\frac12kq_J^2,\\
q_J&=\frac{L_J^2}{\mu kR_e^3}.
\end{align*}

所以转动使势阱最低点向更大键长移动 \(q_J>0\)，同时把能量降低

\[
\Delta E_{\rm cd}
=-\frac{(L_J^2)^2}{2\mu^2kR_e^6}.
\]

因此最低阶转振能谱为

\begin{equation}
\boxed{
E_{vJ}
\simeq
U_e+\hbar\omega_e\left(v+\frac12\right)
+B_eJ(J+1)
-D_e[J(J+1)]^2,
}
\label{eq:rovib-centrifugal-spectrum}
\end{equation}

其中离心畸变常数

\begin{equation}
\boxed{
D_e
=\frac{\hbar^4}{2\mu^2kR_e^6}
=\frac{4B_e^3}{\hbar^2\omega_e^2}.
}
\label{eq:centrifugal-distortion-constant}
\end{equation}

\begin{derivation}[为什么只保留线性 \(q\) 项就能得到最低阶 \(J^4\) 修正]
把 \(1/(R_e+q)^2\) 按 \(q/R_e\) 展开，
\begin{align*}
 \frac1{(R_e+q)^2}
 &=\frac1{R_e^2}\left(1+\frac q{R_e}\right)^{-2}\\
 &=\frac1{R_e^2}
 \left(1-\frac{2q}{R_e}
 +\frac{3q^2}{R_e^2}
 -\frac{4q^3}{R_e^3}+\cdots\right),
\end{align*}
因此有效势中来自转动的三项分别为
\begin{equation*}
 \frac{L_J^2}{2\mu R_e^2}
 -\frac{L_J^2}{\mu R_e^3}q
 +\frac{3L_J^2}{2\mu R_e^4}q^2+\cdots.
\end{equation*}
第一项是常数，正比于 \(L_J^2\sim J(J+1)\)，给出刚性转子能量 \(B_eJ(J+1)\)。第二项线性于 \(q\)，系数正比于 \(L_J^2\)；把它与谐振子势 \(kq^2/2\) 合并并配方，
\begin{equation*}
 \frac12kq^2-\frac{L_J^2}{\mu R_e^3}q
 =\frac12k\left(q-\frac{L_J^2}{\mu kR_e^3}\right)^2
 -\frac12k\left(\frac{L_J^2}{\mu kR_e^3}\right)^2
 =\frac12k(q-q_J)^2
 -\frac{(L_J^2)^2}{2\mu^2kR_e^6},
\end{equation*}
于是新平衡位置
\begin{equation*}
 q_J=\frac{L_J^2}{\mu kR_e^3}\propto L_J^2,
\end{equation*}
最低点相对 \(q=0\) 下降
\begin{equation*}
 \frac12kq_J^2
 =\frac12k\frac{(L_J^2)^2}{\mu^2k^2R_e^6}
 =\frac{(L_J^2)^2}{2\mu^2kR_e^6}\propto (L_J^2)^2\sim [J(J+1)]^2,
\end{equation*}
正好是 \(J^4\) 阶的离心畸变修正 \(-D_e[J(J+1)]^2\)。

二次项 \(3L_J^2q^2/(2\mu R_e^4)\) 在新平衡点 \(q=q_J+\xi\) 附近展开，
\begin{align*}
 \frac{3L_J^2}{2\mu R_e^4}(q_J+\xi)^2
 &=\frac{3L_J^2}{2\mu R_e^4}q_J^2
 +\frac{3L_J^2}{\mu R_e^4}q_J\,\xi
 +\frac{3L_J^2}{2\mu R_e^4}\xi^2.
\end{align*}
其中 \(\xi^2\) 项把振动频率改为 \(\frac12k_{\rm eff}\xi^2\)，\(k_{\rm eff}=k+3L_J^2/(\mu R_e^4)\)，使 \(\omega\) 依赖 \(J\)；\(\xi\) 项进一步微调 \(q_J\)；常数项 \(3L_J^2q_J^2/(2\mu R_e^4)\) 虽然也是 \(J^4\) 阶，但其系数进入 \(v\) 依赖的振转耦合 \(B_v=B_e-\alpha_e(v+1/2)+\cdots\)，而非纯粹的离心畸变常数 \(D_e\)。因此仅用线性项配方已足以得到 \eqnrefs{eq:centrifugal-distortion-constant} 的最低阶 \([J(J+1)]^2\) 修正，二次项属于下一阶的振转精细结构。

\textbf{物理意义。}离心畸变常数 $D_e=4B_e^3/\omega_e^2$ 量化转动分子键拉伸效应：$J$ 越大离心力越强、键越长、转动惯量越大、有效 $B$ 减小。实例：$\mathrm{H}^{35}\mathrm{Cl}$ 的 $B_e=10.59\,\si{\cm^{-1}}$、$\omega_e=2990\,\si{\cm^{-1}}$，给出 $D_e\approx5.3\times10^{-4}\,\si{\cm^{-1}}$，与实验值 $5.28\times10^{-4}\,\si{\cm^{-1}}$ 高度吻合。重分子如 $\mathrm{I}_2$ 的 $D_e\sim10^{-8}\,\si{\cm^{-1}}$，效应极弱。
\end{derivation}

这一步展开有两个独立的小参数：振动本身要求 \(\ell_v/R_e\ll1\)，转动造成的静态位移还要求

\begin{equation}
\frac{q_J}{R_e}
=\frac{\hbar^2J(J+1)}{\mu kR_e^4}\ll1.
\label{eq:centrifugal-control-parameter}
\end{equation}

当 \(J\) 足够大使 \(q_J/R_e\) 不再很小时，简单的 \(B_e-D_eJ^2(J+1)^2\) 截断必须被更高阶 Dunham 展开或直接数值求解 \eqnrefs{eq:diatomic-radial-equation} 取代。

\section{从正常模到振转光谱与 vibronic 跃迁}

本章得到的振动本征态不仅决定红外或拉曼振转谱，也会在电子跃迁中作为核波函数出现。当电子态改变时，势能曲线的平衡位置、曲率乃至正规模式都可能变化，于是两个电子势能面上的核波函数不再彼此正交。它们的重叠正是后文 Franck--Condon 因子的来源。换句话说，Franck--Condon 理论并不是另外引入的一套经验规则，而是本章核 Hamiltonian 在不同电子势能面上的本征态之间的投影问题。

\section{振转光谱的实验数值与典型分子参数}\label{sec:rovib-experiments}

本节给出双原子分子振转光谱的典型实验参数，与\eqnrefs{eq:rovib-centrifugal-spectrum}的理论公式直接对接。

接下来考察双原子分子光谱常数总表。

\begin{center}
\begin{tabular}{lcccccc}
\hline
分子 & \(\tilde\omega_e\) (cm\(^{-1}\)) & \(\tilde B_e\) (cm\(^{-1}\)) & \(\tilde D_e\) (cm\(^{-1}\)) & \(\tilde\alpha_e\) (cm\(^{-1}\)) & \(R_e\) (\AA) & \(\mu\) (amu) \\
\hline
H\(_2\) & 4401.2 & 60.853 & 0.0471 & 3.062 & 0.741 & 0.504 \\
HCl & 2990.9 & 10.593 & 0.00053 & 0.307 & 1.275 & 0.981 \\
N\(_2\) & 2358.6 & 1.998 & 5.8\(\times10^{-6}\) & 0.017 & 1.098 & 7.004 \\
O\(_2\) & 1580.2 & 1.446 & 4.8\(\times10^{-6}\) & 0.016 & 1.208 & 8.000 \\
CO & 2170.0 & 1.931 & 6.5\(\times10^{-6}\) & 0.017 & 1.128 & 6.861 \\
NO & 1904.2 & 1.702 & 5.5\(\times10^{-6}\) & 0.018 & 1.151 & 7.466 \\
I\(_2\) & 214.5 & 0.0374 & 4.2\(\times10^{-9}\) & 0.0004 & 2.667 & 63.45 \\
\hline
\end{tabular}
\end{center}

其中 \(\tilde\omega_e=\omega_e/(2\pi c)\)，\(\tilde B_e=B_e/(hc)\)，\(\tilde D_e=D_e/(hc)\)，\(\tilde\alpha_e\) 是振转耦合常数。约化质量 \(\mu=m_Am_B/(m_A+m_B)\) 以原子质量单位 (amu) 给出。

接下来考察HCl 振转光谱的定量分析。

HCl 是振转光谱教学的典型分子。其红外吸收谱在 \(\tilde\nu\approx2900\,\mathrm{cm}^{-1}\) 附近展现 P 支（\(\Delta J=-1\)）和 R 支（\(\Delta J=+1\)）：

\[
\tilde\nu_R(J)=\tilde\omega_0+2\tilde B_e(J+1)-4\tilde D_e(J+1)^3,\qquad J=0,1,2,\ldots
\]
\[
\tilde\nu_P(J)=\tilde\omega_0-2\tilde B_e J+4\tilde D_e J^3,\qquad J=1,2,3,\ldots
\]

\begin{exbox}[HCl R 支前几条线]
取 \(\tilde\omega_0=2885.9\,\mathrm{cm}^{-1}\)（含零点能修正），\(\tilde B_e=10.593\,\mathrm{cm}^{-1}\)：
\[
\tilde\nu_R(0)=2907.1,\quad\tilde\nu_R(1)=2928.3,\quad\tilde\nu_R(2)=2949.5\,\mathrm{cm}^{-1}
\]
间距 \(\Delta\tilde\nu_R\approx2\tilde B_e=21.19\,\mathrm{cm}^{-1}\)，与实验精确吻合。P 支：\(\tilde\nu_P(1)=2864.7\)，\(\tilde\nu_P(2)=2843.5\,\mathrm{cm}^{-1}\)，间距相同但向低频移动。P 和 R 支之间的缺口（gap）\(\approx4\tilde B_e\approx42.4\,\mathrm{cm}^{-1}\) 是振转光谱的标志性特征。
\end{exbox}

接下来考察非谐性与 Morse 势。

谐振子近似预测等间距能级 \(\Delta E=\hbar\omega_e\)，但实际分子的振动能级间距随 \(v\) 增大而递减。Morse 势
\[
V_M(R)=D_e^{(\mathrm{diss})}\left(1-e^{-a(R-R_e)}\right)^2
\]
给出非谐能级
\[
E_v=\hbar\omega_e\left(v+\frac12\right)-\hbar\omega_ex_e\left(v+\frac12\right)^2
\]
其中非谐常数 \(x_e=\hbar\omega_e/(4D_e^{(\mathrm{diss})})\)，离解能 \(D_e^{(\mathrm{diss})}=\hbar^2a^2/(2\mu)\)。

\begin{center}
\begin{tabular}{lccc}
\hline
分子 & \(\tilde\omega_ex_e\) (cm\(^{-1}\)) & \(D_e^{(\mathrm{diss})}\) (eV) & \(D_0\) (eV) \\
\hline
H\(_2\) & 121.3 & 4.75 & 4.48 \\
HCl & 52.8 & 4.43 & 4.28 \\
N\(_2\) & 14.2 & 9.77 & 9.45 \\
O\(_2\) & 12.1 & 5.12 & 5.02 \\
CO & 13.3 & 11.09 & 10.85 \\
I\(_2\) & 0.6 & 1.54 & 1.51 \\
\hline
\end{tabular}
\end{center}

\begin{exbox}[HCl 的离解能与非谐常数]
HCl 的 \(\tilde\omega_ex_e=52.8\,\mathrm{cm}^{-1}\)，\(\tilde\omega_e=2990.9\,\mathrm{cm}^{-1}\)，给出
\[
D_e^{(\mathrm{diss})}=\frac{\hbar\omega_e}{4x_e}=\frac{\tilde\omega_e^2}{4\tilde\omega_ex_e}\cdot hc
=\frac{2990.9^2}{4\times52.8}\times1.2398\times10^{-4}\approx5.29\,\mathrm{eV}
\]
Morse 势预测值偏高约 19\%，因为 Morse 势忽略了长程 van der Waals 尾巴 \(-C_6/R^6\)，实际离解能 \(D_e^{(\mathrm{diss})}=4.43\,\mathrm{eV}\)。
\end{exbox}

接下来考察振动波函数宽度与零点能。

振动长度 \(\ell_v=\sqrt{\hbar/(\mu\omega_e)}\) 和零点能 \(E_0=\frac12\hbar\omega_e\) 的数值：

\begin{center}
\begin{tabular}{lccc}
\hline
分子 & \(\ell_v\) (\AA) & \(\ell_v/R_e\) & \(E_0\) (eV) \\
\hline
H\(_2\) & 0.168 & 0.227 & 0.273 \\
HCl & 0.113 & 0.089 & 0.185 \\
N\(_2\) & 0.050 & 0.046 & 0.146 \\
O\(_2\) & 0.055 & 0.046 & 0.098 \\
I\(_2\) & 0.077 & 0.029 & 0.013 \\
\hline
\end{tabular}
\end{center}

\(\ell_v/R_e\) 是谐振子近似的控制参数。H\(_2\) 的 \(\ell_v/R_e\approx0.23\) 较大，非谐效应显著；I\(_2\) 的 \(\ell_v/R_e\approx0.03\) 很小，谐振子近似在高振动态仍有效。

\begin{derivation}[Morse 势的精确束缚态能级与因子化代数，即\eqnrefs{eq:morse-exact-levels}]
\emph{策略路线图。}首先写出 Morse 势 Schr\"odinger 方程并作指数变量代换；其次将方程化为合流超几何（Kummer）方程的标准形式；再次由波函数平方可积条件提取量子化约束；然后解出精确能级公式并识别非谐常数；最后验证谐振子极限与离解阈。

\emph{第一步：Morse 势与变量代换。}Morse 势
\begin{equation}
V(r)=D_e\bigl(1-\ee^{-a(r-r_e)}\bigr)^2,
\qquad
D_e>0,\ a>0,\ r\in(-\infty,+\infty),
\label{eq:morse-potential}
\end{equation}
在 \(r\to+\infty\) 时 \(V\to D_e\)（离解阈），在 \(r=r_e\) 时 \(V=0\)（平衡），在 \(r\to-\infty\) 时 \(V\to+\infty\)（硬壁）。径向 Schr\"odinger 方程（\(l=0\)，约化质量 \(\mu\)）为
\begin{equation*}
-\frac{\hbar^2}{2\mu}\frac{\dd^2\chi}{\dd r^2}
+D_e\bigl(1-\ee^{-a(r-r_e)}\bigr)^2\chi
=E\,\chi.
\end{equation*}
首先作无量纲指数代换
\begin{equation}
y=2\lambda\,\ee^{-a(r-r_e)},
\qquad
\lambda\equiv\frac{\sqrt{2\mu D_e}}{\hbar a},
\label{eq:morse-variable-change}
\end{equation}
则 \(y\in(0,\infty)\)，且 \(\dd/\dd r=-ay\,\dd/\dd y\)。方程化为
\begin{equation*}
\frac{\dd^2\chi}{\dd y^2}
+\frac{1}{y}\frac{\dd\chi}{\dd y}
+\frac{1}{y^2}\!\left[-\lambda^2+\lambda y-\frac{\mu(E-D_e)}{\hbar^2a^2}\right]\chi=0.
\end{equation*}

\emph{第二步：化为 Kummer 方程。}其次定义能量参数
\begin{equation*}
\varepsilon\equiv\frac{\sqrt{-2\mu(E-D_e)}}{\hbar a}
=\frac{\sqrt{2\mu(D_e-E)}}{\hbar a},
\end{equation*}
对束缚态 \(E<D_e\) 故 \(\varepsilon\) 为实正数。设 \(\chi(y)=y^\varepsilon\,\ee^{-y/2}\,w(y)\)，代入方程并化简，\(w\) 满足
\begin{equation}
y\,w''+(2\varepsilon+1-y)\,w'+(\lambda-\varepsilon-\tfrac12)\,w=0.
\label{eq:morse-kummer}
\end{equation}
此即合流超几何方程 \(y\,w''+(b-y)\,w'-c\,w=0\)，参数
\begin{equation*}
b=2\varepsilon+1,
\qquad
c=\lambda-\varepsilon-\tfrac12.
\end{equation*}
通解为 \(w=A\,{}_1F_1(c;b;y)+B\,U(c,b,y)\)，其中 \({}_1F_1\) 为 Kummer 函数，\(U\) 为 Tricomi 函数。

\emph{第三步：量子化条件。}再次分析边界条件。当 \(y\to\infty\)（即 \(r\to-\infty\)），Kummer 函数 \({}_1F_1(c;b;y)\sim \ee^y y^{c-b}/\Gamma(c)\) 指数增长，使 \(\chi\sim y^{\varepsilon+c-b}\ee^{y/2}\) 发散。为使波函数平方可积，必须截断级数——即 \(c=-n_v\) 为非正整数：
\begin{equation}
\lambda-\varepsilon-\frac12=-n_v,
\qquad
n_v=0,1,2,\ldots,
\label{eq:morse-quantization}
\end{equation}
等价地 \(\varepsilon=\lambda-n_v-1/2\)。同时要求 \(\varepsilon>0\)（否则 \(y^\varepsilon\) 在 \(y\to0^+\) 不可积），给出束缚态数上限
\begin{equation*}
n_v<\lambda-\frac12,
\qquad
N_{\rm bound}=\lfloor\lambda-\tfrac12\rfloor+1.
\end{equation*}

\emph{第四步：精确能级。}由 \(\varepsilon=\lambda-n_v-1/2\) 与 \(\varepsilon\) 的定义反解 \(E\)：
\begin{equation}
E_{n_v}
=D_e-\frac{\hbar^2a^2}{2\mu}\!\left(\lambda-n_v-\frac12\right)^2
=D_e-\frac{\hbar^2a^2}{2\mu}\!\left(\lambda^2-2\lambda(n_v+\tfrac12)+(n_v+\tfrac12)^2\right).
\label{eq:morse-exact-levels}
\end{equation}
代入 \(\lambda^2=2\mu D_e/(\hbar^2a^2)\)，首项 \(D_e-(\hbar^2a^2/2\mu)\lambda^2=D_e-D_e=0\) 恰消去，整理得
\begin{equation}
\boxed{
E_{n_v}
=\hbar a\sqrt{\frac{2D_e}{\mu}}\!\left(n_v+\frac12\right)
-\frac{\hbar^2a^2}{2\mu}\!\left(n_v+\frac12\right)^2
}
\label{eq:morse-exact-levels-expanded}
\end{equation}
即\eqnrefs{eq:morse-exact-levels}。识别谐振频率与非谐常数
\begin{equation}
\omega_e=a\sqrt{\frac{2D_e}{\mu}},
\qquad
\omega_e x_e=\frac{\hbar a^2}{2\mu},
\qquad
x_e=\frac{\hbar a}{2\sqrt{2\mu D_e}}=\frac{1}{2\lambda},
\label{eq:morse-anharmonic-constants}
\end{equation}
能级写成标准光谱学形式
\begin{equation*}
E_{n_v}=\hbar\omega_e(n_v+\tfrac12)-\hbar\omega_e x_e(n_v+\tfrac12)^2.
\end{equation*}

\emph{第五步：极限检验。}最后验证两个极限。当 \(D_e\to\infty\)（保持 \(\omega_e=a\sqrt{2D_e/\mu}\) 固定，即 \(a\to0\)）时，\(x_e=\hbar\omega_e/(4D_e)\to0\)，\eqnrefs{eq:morse-exact-levels-expanded} 退化为谐振子 \(E_{n_v}=\hbar\omega_e(n_v+1/2)\)。最高束缚态 \(n_v^{\rm max}=\lfloor\lambda-1/2\rfloor\) 的能量
\begin{equation*}
E_{n_v^{\rm max}}\approx D_e-\frac{\hbar^2a^2}{2\mu}\cdot\frac14
=D_e-\frac{\hbar\omega_e}{4}\cdot\frac{\hbar\omega_e}{4D_e}
\approx D_e\quad(\lambda\gg1),
\end{equation*}
恰逼近离解阈 \(D_e\)。离解能 \(D_0=D_e-E_0=D_e-\hbar\omega_e/2+\hbar\omega_ex_e/4\)，零点能修正 \(\hbar\omega_ex_e/4=\hbar^2a^2/(8\mu)\) 来自 Morse 势的非谐性，对 HCl 给出 \(D_0\approx D_e-0.185\,\si{eV}\)，与实验值 \(D_0=4.43\,\si{eV}\) 吻合。


\emph{[Dunham 展开与分子常数的精确确定]} 双原子分子振转能级的 Dunham 展开把所有光谱常数统一在一个幂级数中，是分子光谱学的系统框架。分四步建立。

\emph{第一步：振转能级的 Dunham 展开。}能级 \(E_{v,J}=G(v)+F_v(J)\) 分离振动和转动。Dunham (1932) 给出统一展开：
\begin{equation}
  E_{v,J}=\sum_{k,l}Y_{kl}(v+\tfrac{1}{2})^k[J(J+1)]^l,
  \label{eq:phys-dv-dunham}
\end{equation}
其中 \(Y_{kl}\) 是 Dunham 系数。\(Y_{10}\approx\omega_e\)（谐振频率），\(Y_{20}\approx-\omega_ex_e\)（非谐性），\(Y_{01}\approx B_e\)（转动常数），\(Y_{11}\approx-\alpha_e\)（振转耦合），\(Y_{02}\approx-D_e\)（离心畸变）。Dunham 展开把所有常数统一在 \(\{Y_{kl}\}\) 中，避免了传统光谱学的参数冗余。

\emph{第二步：WKB 推导。}Dunham 展开可从径向 Schrödinger 方程的 WKB 量子化条件系统推导。有效势 \(V_{\rm eff}(r)=V(r)+\frac{\hbar^2J(J+1)}{2\mu r^2}\) 的 WKB 条件 \(\int_{r_1}^{r_2}\sqrt{2\mu(E-V_{\rm eff})}\dd r=(v+\frac{1}{2})\pi\hbar\) 展开为 \(E\) 的幂级数，系数由 \(V(r)\) 的 Taylor 展开确定。Dunham 用展开 \(V(r)=a_0\xi^2(1+a_1\xi+a_2\xi^2+\cdots)\)（\(\xi=r-r_e\)）给出 \(Y_{kl}\) 与 \(\{a_n\}\) 的显式关系。

\emph{第三步：势能重建与 Rydberg--Klein--Rees (RKR)。}Dunham 系数 \(\{Y_{kl}\}\) 反向确定势能 \(V(r)\)：RKR 方法利用 WKB 量子化的逆变换，从 \(G(v)\) 和 \(B_v\)（振动和转动常数）精确重建 \(V(r)\) 的内外转折点。RKR 不假设势能形式（如 Morse 势），从光谱数据直接给出势能曲线——这是分子物理中从光谱到势能的标准方法。

\emph{第四步：高阶修正与超精细。}Dunham 展开的高阶项 \(Y_{30},Y_{21},\ldots\) 对应高阶非谐性和振转耦合，在精密光谱中可测量。超精细结构（核自旋 \(I\) 与电子角动量耦合）给出额外展开 \(E_{\rm hfs}=A_v\boldsymbol I\cdot\boldsymbol J+B_v\frac{3(\boldsymbol I\cdot\boldsymbol J)^2+\frac{3}{2}\boldsymbol I\cdot\boldsymbol J-I(I+1)J(J+1)}{2I(2I-1)J(2J-1)}+\cdots\)（\(A_v\) 磁偶极、\(B_v\) 电四极）。Dunham 展开加超精细展开把双原子分子能级完全参数化，是分子常数测量的标准框架——所有常数从光谱拟合 \(\{Y_{kl},A_v,B_v,\ldots\}\) 获得，物理意义清晰。
\end{derivation}
